\documentclass[11pt,a4paper]{article}
\usepackage[T1]{fontenc}
\usepackage[utf8]{inputenc}
\usepackage{lmodern}
\usepackage[margin=25.5mm,headheight=15pt]{geometry}
\usepackage{amsmath,amssymb,amsthm,mathtools}
\usepackage{booktabs,tabularx,array,microtype,enumitem,xcolor}
\usepackage{fancyhdr,listings,xurl,titlesec}
\definecolor{navy}{RGB}{30,57,82}
\definecolor{ink}{RGB}{32,38,44}
\definecolor{pale}{RGB}{239,243,247}
\usepackage[colorlinks=true,linkcolor=ink,citecolor=ink,urlcolor=navy,bookmarksnumbered=true]{hyperref}
\usepackage{bookmark}
\hypersetup{pdftitle={Sparse Perfect Families in a Single Coarse Class},pdfauthor={Research report prepared for Vladimir Reshetnikov},pdfsubject={A counterexample to C1, quantitative minimal pairs, and representative spectra}}
\urlstyle{same}
\titleformat{\section}{\large\bfseries\color{navy}}{\thesection}{0.7em}{}
\titleformat{\subsection}{\normalsize\bfseries\color{navy}}{\thesubsection}{0.65em}{}
\setlength{\parskip}{3.5pt}
\setlength{\parindent}{1.2em}
\setlength{\emergencystretch}{3em}
\setlist{nosep,leftmargin=*}
\setcounter{tocdepth}{1}
\pagestyle{fancy}\fancyhf{}
\fancyhead[L]{\small\color{navy}Sparse perfect families and coarse classes}
\fancyhead[R]{\small September 2026}
\fancyfoot[C]{\thepage}
\renewcommand{\headrulewidth}{0.3pt}
\newcommand{\N}{\mathbb N}
\newcommand{\Cantor}{2^{\N}}
\newcommand{\Baire}{\N^{\N}}
\newcommand{\strings}{2^{<\N}}
\newcommand{\leT}{\leq_{\mathrm T}}
\newcommand{\eqT}{\equiv_{\mathrm T}}
\newcommand{\degT}{\operatorname{deg}_{\mathrm T}}
\newcommand{\Spec}{\operatorname{Spec}_{\mathrm T}}
\newcommand{\Up}{\operatorname{Up}}
\newcommand{\CD}{\operatorname{CD}}
\newcommand{\Var}{\operatorname{Var}}
\newcommand{\Ham}{\operatorname{Ham}}
\newcommand{\nc}{\mathrm{nc}}
\newcommand{\uc}{\mathrm{uc}}
\newcommand{\zero}{\mathbf 0}
\newcommand{\emptystr}{\varepsilon}
\newcommand{\restr}{\mathbin{\upharpoonright}}
\newcommand{\concat}{\mathbin{{}^{\frown}}}
\newcommand{\code}[1]{\texttt{#1}}
\newtheorem{theorem}{Theorem}[section]
\newtheorem{proposition}[theorem]{Proposition}
\newtheorem{lemma}[theorem]{Lemma}
\newtheorem{corollary}[theorem]{Corollary}
\theoremstyle{definition}\newtheorem{definition}[theorem]{Definition}
\theoremstyle{remark}\newtheorem{remark}[theorem]{Remark}
\lstset{basicstyle=\small\ttfamily,columns=fullflexible,breaklines=true,
  frame=single,rulecolor=\color{navy},showstringspaces=false,
  backgroundcolor=\color{pale},xleftmargin=3pt,xrightmargin=3pt}
\pdfstringdefDisableCommands{\def\N{N}\def\leT{<=T}\def\zero{0}\def\nc{nc}\def\uc{uc}}

\begin{document}
\begin{titlepage}
\vspace*{12mm}
{\small\sffamily\bfseries\color{navy}COMPUTABILITY THEORY \quad / \quad CONSTRUCTION AND SOURCE AUDIT\par}
\vspace{15mm}
{\Huge\bfseries\color{navy}Sparse Perfect Families\par}
\vspace{4mm}
{\LARGE in a Single Coarse Class\par}
\vspace{10mm}
{\large A counterexample to C1, quantitative Turing minimal pairs,\par
and a description of representative spectra\par}
\vspace{13mm}
\noindent\rule{\textwidth}{0.5pt}
\vspace{5mm}
\noindent Prepared for Vladimir Reshetnikov\\[3pt]
18 September 2026
\vspace{9mm}
\begin{abstract}
We attack C1 in the supplied report: must a nonuniform coarse-equivalence class
contain a representative of least Turing degree? The answer is negative. An
important source audit shows that this conclusion already follows from a 2016
theorem of Hirschfeldt, Jockusch, Kuyper, and Schupp, although the coarse clause
is posed again in the displayed Question~7 of Gerdes's 2025 preprint.
We therefore do not claim a new resolution of that question.

The main work here is an independent finite-extension construction with a much
more specific conclusion. For every computable, nondecreasing, unbounded
sublinear function $h$, we construct a perfect family of individually
$1$-generic sets, any two distinct members of which form a Turing minimal pair.
All members agree outside one set $V$ satisfying
$|V\cap[0,n)|\le h(n)$ for every $n$. Two witnesses and the tree labeling are
computable in $0''$. The proof combines a one-bit connectivity lemma with a
global variation budget and perfect-tree fusion. We also identify the Turing
spectra of density-equivalence, uniform coarse-equivalence, and nonuniform
coarse-equivalence classes, and characterize when such spectra have a least
degree. All mathematical proofs are included; finite verification code is
provided separately. Priority for the quantitative refinement is not established.
\end{abstract}
\vfill
\noindent\small The infinite construction is a conventional mathematical proof, not a
claimed Lean certificate. The executable checks verify finite combinatorial
invariants and do not implement a halting-problem oracle.
\end{titlepage}

{\small\tableofcontents}
\newpage

\section{The selected target and the principal result}
\label{sec:result}

\subsection{What C1 asks}
For a total function $f:\N\to\N$, a \emph{coarse description} is a total
function that differs from $f$ on a set of asymptotic density zero. Write
$f\leq_{\nc}g$ when every coarse description of $g$ computes a coarse
description of $f$. The reduction is nonuniform: different descriptions may
require different oracle programs. The selected question is exactly
\begin{equation}\label{eq:C1}
 \forall f\in\Baire\;\exists g\equiv_{\nc}f\;
 \forall k\equiv_{\nc}f\quad g\leT k.
\end{equation}
The supplied report calls this C1 and recommends a noncomputable coarse class
containing a Turing minimal pair as a sufficient counterexample certificate
\cite[Section~5, C1 and its diagnostic propositions]{Input}.

A \emph{Turing minimal pair} here means two noncomputable sets $A,B$ such that
every total numerical function computable from both is computable. This is a
statement about their Turing degrees, not their coarse degrees. In our
construction their coarse degrees will be \emph{equal}, rather than forming a
minimal pair in the coarse degrees.

\subsection{The quantitative theorem}
\begin{definition}[Sublinear order]
A function $h:\N\to\N$ is a sublinear order if it is computable,
nondecreasing, unbounded, and $h(n)/n\to0$ as $n\to\infty$ through positive
integers. The value $h(0)$ may be zero.
\end{definition}

\begin{theorem}[Sparse perfect family]\label{thm:main}
For every sublinear order $h$, there are a nonempty perfect set
$P\subseteq\Cantor$ and a set $V\subseteq\N$ with the following properties.
\begin{enumerate}[label=\textup{(\roman*)}]
\item Every $X\in P$ is $1$-generic.
\item If $X,Y\in P$ are distinct, then every total function $f$ with
$f\leT X$ and $f\leT Y$ is computable.
\item All members of $P$ agree outside $V$, and
\begin{equation}\label{eq:budget-main}
                 |V\cap[0,n)|\le h(n)\qquad(n\in\N).
\end{equation}
\item A compatible labeling of a full binary tree whose paths are $P$ is
computable in $0''$. In particular, $P$ contains two distinct members
$A,B\leT0''$.
\end{enumerate}
Consequently every two distinct members form a Turing minimal pair, and all
members belong to one nonzero uniform coarse class and one nonzero nonuniform
coarse class. Neither class has a least Turing degree.
\end{theorem}

The argument is developed in Sections~\ref{sec:extension}--\ref{sec:perfect}.
The construction uses no minimal-pair theorem or cone-avoidance theorem as a
black box. Its noncomputability requirement is obtained by meeting the
ordinary $1$-genericity requirements.

An explicit permissible bound is
\begin{equation}\label{eq:log-bound}
                     h(n)=\lfloor\log_2(n+1)\rfloor.
\end{equation}
Thus the same coarse class contains continuum many pairwise Turing-minimal-pair
representatives, even though, among their first million bits, all variation
across the \emph{entire family} is confined to at most $19$ coordinates.
The theorem does not say that those coordinates can be located computably.
For example, the bound in \eqref{eq:log-bound} also permits the disagreement
positions to be much farther apart than exponential spacing.

\subsection{What is established, and what is not a novelty claim}
The two main conclusions have different provenance. The negative answer to
\eqref{eq:C1} is an implication of previously published work; the exact bridge
is supplied in Section~\ref{sec:audit}. The independently developed work in
this report is the quantitative sparse-family construction and its supporting
lemmas. A targeted literature search did not identify this exact combination
of a perfect family, pairwise Turing minimality, one common variation set, and
an arbitrary prescribed computable sublinear bound. That search is not a
proof of priority. The quantitative theorem should therefore be regarded as
a detailed research result offered for mathematical review, not advertised as
a certified first occurrence in the literature.

The distinction matters for the user's original objective. We obtain a full
counterexample to the selected formula, not merely a conditional obstruction.
But the source audit prevents us from describing it as a newly solved
literature-open problem.

\section{Definitions and the source-status correction}
\label{sec:audit}

\subsection{Density, descriptions, and the two reductions}
For $S\subseteq\N$ and $n>0$, put
\[
 \rho_n(S)=\frac{|S\cap[0,n)|}{n}.
\]
For total $f,g$, write $f=_{\rho}g$ if
$\rho_n(\{i:f(i)\ne g(i)\})\to0$, and put
\[
                     \CD(f)=\{D:D=_{\rho}f\}.
\]
A finite union of density-zero sets has density zero: at each prefix, the
cardinality of a union is at most the sum of the cardinalities. It follows
that $=_{\rho}$ is an equivalence relation and
\begin{equation}\label{eq:identical-CD}
                f=_{\rho}g\quad\Longrightarrow\quad\CD(f)=\CD(g).
\end{equation}

The definitions used here, as in C1, are
\begin{align*}
 f\leq_{\nc}g
 &\quad\Longleftrightarrow\quad
   \forall D\in\CD(g)\;\exists E\in\CD(f)\quad E\leT D,\\
 f\leq_{\uc}g
 &\quad\Longleftrightarrow\quad
   \exists\Phi\;\forall D\in\CD(g)\quad\Phi^D\in\CD(f).
\end{align*}
The second definition includes totality of $\Phi^D$ on each indicated input
oracle. Identity is a uniform reduction in both directions in
\eqref{eq:identical-CD}. Hence
\begin{equation}\label{eq:equiv-inclusions}
          f=_{\rho}g\quad\Longrightarrow\quad
          f\equiv_{\uc}g\quad\Longrightarrow\quad f\equiv_{\nc}g.
\end{equation}
These conventions agree with the explicit nonuniform formula in Gerdes's
Question~7 and with the coarse-reducibility definitions in the earlier
literature \cite{Gerdes,HJKS}.

For binary $f$, allowing natural-valued descriptions causes no issue here.
The computable clipping map $c(0)=0$ and $c(m)=1$ for $m>0$ fixes both binary
values, so
\[
 \{i:c(D(i))\ne f(i)\}\subseteq\{i:D(i)\ne f(i)\}.
\]
A natural-valued computable coarse description would therefore yield a binary
computable coarse description.

\begin{lemma}[Generic sets are not coarsely computable]\label{lem:generic}
Every $1$-generic binary set has no computable coarse description, even among
natural-valued functions.
\end{lemma}
\begin{proof}
Recall that $X$ is $1$-generic if, for each c.e. set $W\subseteq\strings$,
either some prefix of $X$ belongs to $W$, or some prefix of $X$ has no
extension in $W$.
Fix a total computable $C:\N\to\N$, a rational $q<1$ with $q\ge0$, and an
integer $N$. The following set of finite binary strings is computable:
\[
 W_{C,q,N}=\left\{\sigma:\ |\sigma|>N,
 \quad\frac{|\{i<|\sigma|:\sigma(i)\ne C(i)\}|}{|\sigma|}>q\right\}.
\]
It is dense. Given any prefix, extend it by sufficiently many bits opposite
to $C$; at position $i$ use $1$ if $C(i)=0$ and $0$ otherwise. All appended
bits are disagreements, so the disagreement fraction eventually exceeds $q$.
A $1$-generic cannot avoid a dense set of strings, and therefore has a prefix
in $W_{C,q,N}$. Varying $N$ shows that the upper density of its disagreements
with $C$ is at least $q$. Varying rational $q<1$ makes this upper density equal
to $1$, which rules out a coarse description.
\end{proof}

\begin{lemma}[Minimal-pair certificate]\label{lem:certificate}
Suppose $A=_{\rho}B$, $A$ has no computable coarse description, and every total
function below both $A$ and $B$ is computable. Then neither $[A]_{\nc}$ nor
$[A]_{\uc}$ has a representative of least Turing degree.
\end{lemma}
\begin{proof}
First, $[A]_{\nc}$ contains no computable member. If a computable $g$ belonged
to it, the relation $A\leq_{\nc}g$, applied to the coarse description $g$ of
itself, would produce a computable coarse description of $A$.
Both classes contain $A,B$, by \eqref{eq:equiv-inclusions}. A least-degree
representative $k$ of either class would satisfy $k\leT A,B$ and so would be
computable, contrary to the preceding observation. This argument allows $k$
to be natural-valued and proves absence of a \emph{least} degree, not absence
of all minimal elements.
\end{proof}

\subsection{Why the negative answer already follows from the literature}
Gerdes's arXiv version dated 9 August 2025 explicitly asks
\eqref{eq:C1} in Question~7, on printed page~64 \cite{Gerdes}.
In contrast, Hirschfeldt--Jockusch--Kuyper--Schupp define the information common
to all coarse descriptions by
\[
 X^{\mathfrak c}=
 \{C\subseteq\N:\ \forall D\in\CD(X)\ (D\text{ binary}\Rightarrow C\leT D)\}.
\]
Their Theorem~4.2 states that if $X$ is $1$-generic, then $X^{\mathfrak c}$
consists only of computable sets \cite[Theorem~4.2]{HJKS}.
The paper was published in the \emph{Journal of Symbolic Logic} in 2016;
the inspected preprint is dated May 2015 \cite{HJKS}.

\begin{proposition}[The older theorem refutes C1]\label{prop:old}
For every $1$-generic $X$, neither $[X]_{\nc}$ nor $[X]_{\uc}$ has a least
Turing degree.
\end{proposition}
\begin{proof}
Assume $g$ is a least-degree representative of one of these classes.
Every binary coarse description $D$ of $X$ satisfies $D=_{\rho}X$, so is a
member of both classes. Thus $g\leT D$ for every such $D$.
The graph of the total function $g$, coded as a set of natural numbers, is
therefore in $X^{\mathfrak c}$. By the cited theorem it is computable, hence
$g$ itself is computable: on input $n$, search the computable graph for the
unique pair $(n,m)$. But $X\leq_{\nc}g$ now yields a computable coarse
description of $X$, contradicting Lemma~\ref{lem:generic}.
\end{proof}

This is an exact quantifier match, not an analogy between distinct degree
structures. The logical step is that \emph{every coarse description is itself
a representative of the same coarse degree}. No minimal-pair construction is
needed for Proposition~\ref{prop:old}. The rest of this report goes beyond this
literature bridge by constructing exceptionally close minimal-pair witnesses
directly.

\section{The Turing spectrum of a coarse class}
\label{sec:spectra}

For a class $\mathcal C$ of total functions, write
\[
 \Spec(\mathcal C)=\{\degT(g):g\in\mathcal C\},\qquad
 \Up(\mathcal C)=\{\mathbf b:\exists g\in\mathcal C\ (\degT(g)\le\mathbf b)\}.
\]
The following elementary observation strengthens the upward-closure diagnostic
in the supplied report \cite[Section~5]{Input}.

\begin{theorem}[Equality of representative spectra]\label{thm:spectra}
For every total $f:\N\to\N$,
\begin{equation}\label{eq:spectra}
 \Spec([f]_{\rho})
   =\Spec([f]_{\uc})
   =\Spec([f]_{\nc})
   =\Up(\CD(f)).
\end{equation}
If $f$ is binary, the same equalities hold when all representative classes and
descriptions in the display are restricted to binary functions. Their degree
spectra also coincide with the unrestricted spectra in \eqref{eq:spectra}.
\end{theorem}
\begin{proof}
The first two inclusions follow from \eqref{eq:equiv-inclusions}.
If $g\equiv_{\nc}f$, then applying $f\leq_{\nc}g$ to the description $g$
itself produces $D\in\CD(f)$ with $D\leT g$. Consequently
$\Spec([f]_{\nc})\subseteq\Up(\CD(f))$.

For the reverse inclusion, take a degree $\mathbf b$ computing
$D\in\CD(f)$, and a binary set $B$ of degree $\mathbf b$. Define
\[
 Y(k)=\begin{cases}
        B(n),&k=2^n\text{ for some }n,\\
        D(k),&\text{otherwise}.
       \end{cases}
\]
The powers of two have density zero, so $Y=_{\rho}D=_{\rho}f$.
Moreover $Y\leT B$ because $D\leT B$, and $B\leT Y$ by reading the
coordinates $2^n$. Hence $\degT(Y)=\mathbf b$, proving the reverse inclusion.
For binary $f$, first clip any natural-valued $D$ to a binary description
without increasing its degree, then perform the same coding. This gives a
binary representative of every degree in the unrestricted spectrum.
\end{proof}

These spectra are always upward closed. Uniform and nonuniform equivalence
need not be identical relations, but their classes \emph{at a fixed $f$} have
identical sets of Turing degrees. In particular, least-degree existence is
the same for $[f]_{\rho}$, $[f]_{\uc}$, and $[f]_{\nc}$.

\subsection{A least-degree criterion by robust block coding}
For a binary $B$, define $J(B)$ by
\begin{equation}\label{eq:J}
 J(B)(0)=0,\qquad J(B)(m)=B(n)\quad(2^n\le m<2^{n+1}).
\end{equation}
Clearly $J(B)\eqT B$. This block code is used below only with nonuniform
reductions; the finite correction in its decoder must not be silently
promoted to a uniform operation.

\begin{lemma}[Block recovery]\label{lem:block}
Every coarse description $D$ of $J(B)$ computes $B$.
\end{lemma}
\begin{proof}
Clip $D$ to a binary function first. On the block
$I_n=[2^n,2^{n+1})$, let $B^*(n)$ be the majority value of $D$, breaking ties
in favor of zero. If $E$ is its error set relative to $J(B)$, then
\[
 \frac{|E\cap I_n|}{2^n}
 \le 2\rho_{2^{n+1}}(E)\longrightarrow0.
\]
For all sufficiently large $n$, fewer than half the block positions are
errors, so $B^*(n)=B(n)$. Thus $B^*\leT D$ and $B=^*B^*$, which gives
$B\leT D$ by hard-coding the finitely many exceptional bits. Neither their
number nor their locations are asserted to be computable uniformly from $D$.
\end{proof}

\begin{theorem}[When a least degree exists]\label{thm:least-characterization}
For every total $f$, the following are equivalent:
\begin{enumerate}[label=\textup{(\roman*)}]
\item $[f]_{\nc}$ contains a least Turing degree.
\item There is a binary $B$ with $f\equiv_{\nc}J(B)$.
\end{enumerate}
In this case the least degree is $\degT(B)$. By Theorem~\ref{thm:spectra}, it is
also the least degree of $[f]_{\uc}$ and of $[f]_{\rho}$.
\end{theorem}
\begin{proof}
Suppose $g\equiv_{\nc}f$ has least degree, and choose $B\eqT g$ binary.
Every $D\in\CD(f)$ is a representative of $[f]_{\nc}$, so computes $g$ and
therefore computes $J(B)$. This proves $J(B)\leq_{\nc}f$.
Also $g$, and hence $B$, computes a coarse description of $f$. By
Lemma~\ref{lem:block}, every coarse description of $J(B)$ computes $B$, hence
computes a coarse description of $f$. Thus $f\leq_{\nc}J(B)$.

Conversely, suppose $f\equiv_{\nc}J(B)$. The latter is a representative of
degree $\degT(B)$. If $g\equiv_{\nc}f$, applying $J(B)\leq_{\nc}g$ to $g$
itself yields a coarse description of $J(B)$ computable from $g$.
Lemma~\ref{lem:block} gives $B\leT g$. So $\degT(B)$ is least.
\end{proof}

The criterion can also be phrased without a code: a least degree exists
precisely when some degree computed by \emph{every} coarse description also
computes \emph{one} coarse description. Merely possessing common information
is not enough. The code $J$ packages both requirements into a coarse degree.

\section{A finite-extension lemma with a one-bit reserve}
\label{sec:extension}

\subsection{Finite oracle computations}
Let $(\Phi_e)_{e\in\N}$ enumerate the partial oracle functionals with numerical
outputs. The assertion $\Phi_e^\sigma(n)\downarrow=a$ means that a finite
computation halts with value $a$ using only oracle coordinates below
$|\sigma|$. Its value persists under every extension of $\sigma$.
A missing finite convergence witness is not a proof of divergence on all
infinite extensions.

We write $\sigma u$ for concatenation, and require paired suffixes $u,v$ to
have equal length. Their Hamming distance is
\[
          \Ham(u,v)=|\{i<|u|:u(i)\ne v(i)\}|.
\]
The next lemma isolates the main combinatorial mechanism independently of
density or of a particular budget.

\begin{lemma}[One-bit connectivity trilemma]\label{lem:trilemma}
Fix finite binary strings $\sigma,\tau$, and a collection $\mathcal A$ of pairs
of equal-length binary suffixes. Assume
\begin{equation}\label{eq:one-bit}
             \Ham(u,v)\le1\quad\Longrightarrow\quad(u,v)\in\mathcal A.
\end{equation}
For any two functionals $\Phi,\Psi$, at least one of the following holds:
\begin{description}[style=nextline]
\item[D: admissible disagreement.]
There are $n$ and $(u,v)\in\mathcal A$ such that both
$\Phi^{\sigma u}(n)$ and $\Psi^{\tau v}(n)$ converge, to different values.
\item[P: a partiality cylinder.]
For some $n,u$, either
\begin{equation}\label{eq:partial-cylinder}
        \forall w\in\strings\quad\Phi^{\sigma u w}(n)\uparrow,
\end{equation}
or the corresponding assertion holds for $\Psi$ above $\tau u$.
\item[C: a computable common output.]
There is a total computable $b:\N\to\N$ such that, for every $n,u$,
\begin{align*}
 \Phi^{\sigma u}(n)\downarrow&\quad\Longrightarrow\quad
      \Phi^{\sigma u}(n)=b(n),\\
 \Psi^{\tau u}(n)\downarrow&\quad\Longrightarrow\quad
      \Psi^{\tau u}(n)=b(n).
\end{align*}
\end{description}
In C, the conclusion is valid for \emph{all} finite suffixes, not merely
admissibly paired ones. Hence it remains true after arbitrary subsequent
extensions on either side.
\end{lemma}

\begin{proof}
Assume D and P both fail. Fix an input $n$.
Failure of P says that each functional can be made convergent above
\emph{every} finite suffix on its side. More explicitly, for each $u$ there
is $w$ with $\Phi^{\sigma u w}(n)\downarrow$, and similarly for $\Psi$.
These convergence conditions are dense open conditions on appended suffixes.

Take any two finite convergence witnesses, on either of the two sides, with
suffixes $u,v$. Pad the shorter suffix with arbitrary bits until both have
the same length $d$. Padding preserves the witnessed computations.
In the $d$-dimensional binary cube, join $u$ to $v$ by a path
\[
             u=u_0,u_1,\ldots,u_m=v
\]
that flips one differing coordinate at a time. Consecutive vertices have
Hamming distance one.

There is a \emph{single} common suffix $w$ such that all the finitely many
computations
\begin{equation}\label{eq:simultaneous}
  \Phi^{\sigma u_iw}(n),\quad\Psi^{\tau u_iw}(n)\qquad(0\le i\le m)
\end{equation}
converge. To see this without invoking a compactness theorem, process these
$2(m+1)$ requirements one at a time. If a current common suffix is $w_0$,
density supplies an additional suffix for the next required computation.
Append it to $w_0$ on every occurrence. All previously obtained finite
convergences remain true. After finitely many steps we have $w$.

The pairs $(u_iw,u_iw)$ are admissible, so failure of D gives
\begin{equation}\label{eq:diagonal-equality}
      \Phi^{\sigma u_iw}(n)=\Psi^{\tau u_iw}(n).
\end{equation}
Also $(u_iw,u_{i+1}w)$ is admissible: appending the same word leaves its
Hamming distance equal to one. Again by failure of D,
\begin{equation}\label{eq:edge-equality}
      \Phi^{\sigma u_iw}(n)=\Psi^{\tau u_{i+1}w}(n).
\end{equation}
Equations \eqref{eq:diagonal-equality} and \eqref{eq:edge-equality} propagate
one value along the whole path, on both sides. In particular, the two
arbitrary convergence witnesses chosen initially had the same value.
Thus every convergent finite computation on input $n$, on either side, has
one common value.

Finally, to compute $b(n)$, dovetail all finite computations
$\Phi^{\sigma u}(n)$ until one converges and output its value. The search
terminates because P fails already at $u=\emptystr$. The preceding paragraph
proves uniqueness and correctness of the output. The finite string $\sigma$
and the index of $\Phi$ are constants in this ordinary program. Even if
$\sigma$ was chosen using an oracle, hard-coding a \emph{finite} string does
not make the resulting program an oracle computation. This proves C.
\end{proof}

\subsection{Why the partiality case cannot be omitted}
A failed disagreement search is not itself sufficient to invoke C. The proof
uses the absence of a \emph{whole cylinder} of divergence to obtain density
of every convergence condition. A bounded simulation cannot certify this
absence. Conversely, after finding a genuine cylinder as in
\eqref{eq:partial-cylinder}, every infinite extension through that cylinder
has a partial output; no subsequent injury or restraint argument is needed.

The one-bit condition is also essential to this lemma. If only diagonal
pairs $(u,u)$ were allowed, the two identity functionals above identical
prefixes would never disagree on an admissible pair, although they could
compute a noncomputable common infinite tail. Diagonal agreement alone does
not force a computable output. Connectivity through one-bit edges removes
exactly this obstruction.

\section{The global variation budget}
\label{sec:budget}

\subsection{Finite frontiers and buffering}
A \emph{frontier} is a finite nonempty indexed family
$F=(\sigma_i)_{i<k}$ of binary strings of a common length $L$. Put
\[
 \Var(F)=\{n<L:\exists i,j<k\ (\sigma_i(n)\ne\sigma_j(n))\}.
\]
This is a global variation set, rather than the disagreement set of one
selected pair. Call $F$ \emph{legal} if
\begin{equation}\label{eq:frontier-legal}
          |\Var(F)\cap[0,n)|\le h(n)\qquad(0\le n\le L).
\end{equation}
Let $r=|\Var(F)|$. Appending the same word to every member creates no new
variation. Legality is preserved because $h$ is nondecreasing.

\begin{lemma}[Buffering]\label{lem:buffer}
Every legal frontier can be extended by a common string of zeros to a legal
frontier of length $L^*$ satisfying $h(L^*)\ge r+1$. The search for such
$L^*$ is computable.
\end{lemma}
\begin{proof}
Since $h$ is unbounded and nondecreasing, the search through lengths
$L,L+1,\ldots$ for the indicated inequality terminates. Common zeros add no
variation, and at new prefix lengths the old count $r$ is bounded by $h$.
\end{proof}

After buffering, fix two different frontier members $\sigma_i,\sigma_j$.
A pair of suffixes $(u,v)$ of common length $d$ is \emph{admissible} if
\begin{equation}\label{eq:admissible}
 r+\Ham(u\restr t,v\restr t)\le h(L+t)
                         \qquad(0\le t\le d).
\end{equation}
This is a decidable finite condition. The inequality $h(L)\ge r+1$ and
monotonicity show that \emph{every} diagonal or one-bit pair of suffixes, of
any length, is admissible. Thus Lemma~\ref{lem:trilemma} applies.

\subsection{Lifting a pair extension to every leaf}
\begin{lemma}[Global lifting]\label{lem:lift}
For admissible $(u,v)$, extend $\sigma_j$ by $v$ and every other frontier
member, including $\sigma_i$, by $u$. The resulting frontier $F^*$ is legal.
Its variation set is exactly
\begin{equation}\label{eq:lift-variation}
 \Var(F^*)=\Var(F)\ \cup\ \{L+t:t<d,\ u(t)\ne v(t)\}.
\end{equation}
\end{lemma}
\begin{proof}
Old coordinates are unchanged. At a new coordinate, at least one leaf has
the $u$-value and the selected $j$-leaf has the $v$-value; all leaves have
one of these values. The coordinate varies exactly when those two values
differ. This proves \eqref{eq:lift-variation}. For prefixes up to $L$ use old
legality; for later prefixes use \eqref{eq:admissible}.
\end{proof}

There is a useful asymmetry here: a pairwise disagreement task can affect
many leaves, but it costs only the differing coordinates of two suffixes,
not the number of leaves. In the partiality case P, append the same witness
$u$ to every frontier member. This costs no new variation and places the
selected side permanently inside the chosen partiality cylinder. In case C,
no special extension is necessary.

\begin{proposition}[One permanent pair requirement]\label{prop:one-requirement}
Given a legal frontier, two distinct leaves $i,j$, and indices $e,d$, a finite
legal extension can ensure that, for all infinite extensions $X$ of the new
$i$-leaf and $Y$ of the new $j$-leaf,
\begin{equation}\label{eq:R}
  \Phi_e^X=\Phi_d^Y\text{ is total}
                 \quad\Longrightarrow\quad\Phi_e^X\text{ is computable}.
\end{equation}
The finite extension can be selected using $0''$.
\end{proposition}
\begin{proof}
Buffer and apply Lemma~\ref{lem:trilemma} to the admissible pairs in
\eqref{eq:admissible}. In D use Lemma~\ref{lem:lift} to make the disagreement
permanent. In P append the common witness so one output is permanently
partial. In C all total outputs are already the computable function supplied
by that lemma. Each outcome is preserved by \emph{arbitrary} further
extensions of all leaves.

For effectiveness, D is a $\Sigma^0_1$ question: its witness consists of
finite suffixes, an input, and two finite halting computations, together with
finitely many computable budget checks. Decide it using $0'$ and search for
a witness when it holds. If D fails, existence of P is $\Sigma^0_2$; its
precise form for the first side is
\[
 \exists n,u\;\forall w,t\quad
            \neg\bigl[\Phi_{e,t}^{\sigma_i u w}(n)\downarrow\bigr].
\]
It can be decided by $0''$. If true, find a witness $(n,u)$ using the
$0'$-decidable inner $\Pi^0_1$ property. If both D and P fail, C follows from
the proved lemma. There is no additional oracle decision in case C.
\end{proof}

\section{A two-real construction first}
\label{sec:pair}

Before handling a perfect family, it is helpful to isolate the complete
counterexample with only two paths.

\begin{theorem}[Sparse $0''$-computable minimal pair]\label{thm:pair}
For every sublinear order $h$, there are $1$-generic $A,B\leT0''$ with
\[
 |(A\mathbin\triangle B)\cap[0,n)|\le h(n)\quad(n\in\N)
\]
such that every total function computable from both is computable.
\end{theorem}
\begin{proof}
Maintain two finite strings $\sigma,\tau$ of equal length forming a legal
frontier. Begin with both strings empty. Schedule all requirements
\[
 R_{e,d}:\quad\Phi_e^A=\Phi_d^B\text{ total}
                       \ \Longrightarrow\ \Phi_e^A\text{ computable},
\]
and all ordinary genericity requirements $G_e^A,G_e^B$, using an effective
enumeration of c.e. sets of strings $(W_e)$.

Handle $R_{e,d}$ by Proposition~\ref{prop:one-requirement}, treating the two
strings as distinct indexed leaves even if their current values happen to
coincide. The proposition and its proof do not require incompatible strings;
only two distinct indices are needed for the lifting argument.

For $G_e^A$, ask whether $\sigma$ has an extension $\sigma u\in W_e$.
If so, find one and append the same $u$ to both current strings. Then $A$
meets $W_e$. If not, the current $\sigma$ already witnesses avoidance, and
no extension is required. Handle $G_e^B$ in the same way using $\tau$.
These are $0'$-decidable existence questions. Common extensions preserve
legality. Append a common zero after each scheduled task to make the lengths
tend to infinity.

Let $A,B$ be the unions of their respective strings. Every genericity
requirement is satisfied permanently, so both paths are $1$-generic. Any
total common function has indices $e,d$ and is computable by the permanent
satisfaction of $R_{e,d}$. By Lemma~\ref{lem:generic}, both paths are
noncomputable. They cannot be equal, or even Turing equivalent: their own
noncomputable characteristic function would then be a common total function.

For each $n$, choose a stage at which the current length is at least $n$.
No later stage changes those coordinates, and the legal-frontier invariant
proves the displayed bound. Since $h(n)/n\to0$, we have $A=_{\rho}B$.
All stage choices are computable in $0''$ by the preceding analysis, so
$A,B\leT0''$.
\end{proof}

\begin{corollary}[A direct negative answer to C1]\label{cor:C1}
The formula \eqref{eq:C1} is false, including when its counterexample $f$ is
required to be binary and computable in $0''$. The analogous uniform
coarse-equivalence formula is also false.
\end{corollary}
\begin{proof}
Use $f=A$ from Theorem~\ref{thm:pair} and apply
Lemmas~\ref{lem:generic} and~\ref{lem:certificate}.
\end{proof}

This is a complete construction of a mathematical counterexample. It is not
an ordinary executable generator of its bits: a program that makes the
specified stage decisions needs the indicated oracle. The proof says exactly
which oracle questions suffice.

\section{Perfect-family fusion}
\label{sec:perfect}

We now prove the full Theorem~\ref{thm:main}. Besides strengthening the result,
the perfect-family version exposes why the budget must track variation across
\emph{all} leaves rather than just one pair.

\subsection{The finite-stage construction}
At level $s$, keep a family of strings
\[
             F_s=(T_s(\alpha))_{\alpha\in2^s}
\]
of a common length $L_s$. They are pairwise incompatible, legal, and extend
their previously frozen ancestor labels. Set
$T_0(\emptystr)=\emptystr$.
Construct the next level by the following finite operations.

\paragraph{Split every leaf.}
Buffer all current leaves as in Lemma~\ref{lem:buffer}. If the buffered
strings are $\widehat T_s(\alpha)$, replace each by the two strings
\[
       \widehat T_s(\alpha)0,\qquad\widehat T_s(\alpha)1.
\]
All splits occur at the same coordinate. They add exactly \emph{one} global
variation coordinate, irrespective of the number of leaves. The reserve
$h(L)\ge r+1$ and monotonicity imply legality after the split.
Label the new leaves by $\alpha0,\alpha1$.

\paragraph{Make every leaf generic to the current level.}
For each current leaf $\beta\in2^{s+1}$ and each $e\le s$, ask whether that
leaf's current string has an extension in $W_e$. When it does, choose a
witness suffix and append that suffix to \emph{all} leaves. When it does not,
record avoidance; no extension is needed. This is a finite sequence of
operations. It adds no variation and permanently settles the indicated
genericity requirements for every future descendant of the selected leaf.

\paragraph{Treat every pair to the current level.}
For each ordered pair of distinct current leaves $\beta,\gamma$ and each
$e,d\le s$, perform Proposition~\ref{prop:one-requirement} for their current
strings. Buffer again before each application, because preceding operations
may have consumed the earlier reserve. Lift the chosen extension to the
whole frontier as prescribed there. All previous requirements remain
satisfied, since their certificates persist under arbitrary extensions.

\paragraph{Freeze the next level.}
The resulting strings define $T_{s+1}(\beta)$ and their common length
$L_{s+1}$. They extend the appropriate old labels. Children of one parent
remain incompatible at their splitting coordinate, and descendants of
different old leaves remain incompatible as well. At least one coordinate
has been appended, so $L_{s+1}>L_s$.

This completes the recursive definition. Each level contains only finitely
many tasks, although no practical bound on their witness-search lengths is
asserted. Each successful existential search terminates by its prior oracle
answer; an unsuccessful existential question is answered by the oracle, not
by waiting indefinitely.

\subsection{The set of paths is perfect}
For an address $Z\in\Cantor$, define
\begin{equation}\label{eq:path-map}
                   F(Z)=\bigcup_s T_s(Z\restr s),
              \qquad P=\{F(Z):Z\in\Cantor\}.
\end{equation}
The union is a total binary sequence because labels are nested and lengths
tend to infinity. Agreement of two addresses through level $s$ implies
agreement of their images through $L_s$, so $F$ is continuous. Distinct
addresses eventually choose different children, whose labels are
incompatible; hence $F$ is injective. By compactness of Cantor space, its
image $P$ is closed.

To verify directly that no point of $P$ is isolated, fix $F(Z)$ and a finite
prefix length $n$. Choose $s$ with $L_s\ge n$, keep the address $Z\restr s$,
and change a later address bit. The new image is distinct but agrees with
$F(Z)$ through its first $n$ coordinates. Thus $P$ is a nonempty perfect set,
and its cardinality is $2^{\aleph_0}$.

\subsection{Genericity and minimal pairs}
Fix $X=F(Z)\in P$ and $e$. At every sufficiently late stage, the current leaf
on $Z$ is treated for $W_e$. Its treatment either leaves a prefix in $W_e$ or
leaves a prefix with no extension in $W_e$. All later descendants preserve
that fact. Therefore $X$ is $1$-generic.

Now fix distinct $X=F(Z)$ and $Y=F(W)$ and suppose a total function $b$ is
computed by $\Phi_e^X$ and $\Phi_d^Y$. Choose a level after the addresses
have separated and after the indices $e,d$ enter the finite schedule. At
that level their distinct leaves are treated for precisely these two
functionals. Outcome D would contradict the equality of the outputs, and
outcome P would contradict totality. Thus their common output is computable,
as ensured by outcome C. Every member is noncomputable by
Lemma~\ref{lem:generic}; hence every distinct pair is a Turing minimal pair.

The paths are only asserted to be \emph{individually} $1$-generic. Relative
or mutual genericity is not part of the construction or the proof.

\subsection{The common sparse variation set}
Every old leaf has descendants forever, so every old variation coordinate
remains a variation coordinate at later frontiers. Define
\[
                         V=\bigcup_s\Var(F_s).
\]
No coordinate can begin to vary after its bits have already been frozen
on all current leaves. Consequently, for $L_s\ge n$,
\[
               V\cap[0,n)=\Var(F_s)\cap[0,n).
\]
Legality gives $|V\cap[0,n)|\le h(n)$ for every $n$.
At any coordinate outside $V$, all labels long enough to cover that
coordinate have the same bit; thus all paths agree there. In particular,
for any $X,Y\in P$,
\[
       (X\mathbin\triangle Y)\subseteq V,
       \qquad\rho_n(X\mathbin\triangle Y)\le\frac{h(n)}n\longrightarrow0.
\]
They are all density-equivalent and hence uniformly coarse-equivalent.
Lemmas~\ref{lem:generic} and~\ref{lem:certificate} give the advertised absence
of least Turing degrees.

\subsection{Effective presentation and two arithmetical witnesses}
The finite list $F_s$ is computable uniformly in $s$ from $0''$: splitting and
budgeting are computable, genericity decisions use $0'$, and pair decisions
use at most $0''$. To compute $F(Z)(n)$ from $Z\oplus0''$, construct levels
until $L_s>n$, then read the appropriate label using $Z\restr s$.
In particular, two fixed computable addresses, for example $0^\N$ and
$10^\N$, yield distinct paths $A,B\leT0''$. Also $V\leT0''$, by constructing
one level beyond the queried coordinate and checking its finite variation set.

This proves all parts of Theorem~\ref{thm:main}. It does \emph{not} assert that
every member of the uncountable family is computable in $0''$. The tree
labeling is $0''$-computable; unrestricted addresses need not be.

\section{Consequences and sharp boundary checks}
\label{sec:boundaries}

\subsection{Unboundedness is the exact threshold for monotone budgets}
\begin{corollary}[Budget dichotomy]\label{cor:dichotomy}
Let $h:\N\to\N$ be computable and nondecreasing. There are noncomputable
$A,B$ forming a Turing minimal pair with
\[
             |(A\mathbin\triangle B)\cap[0,n)|\le h(n)
                         \quad\text{for every }n
\]
if and only if $h$ is unbounded. In the unbounded case, one can additionally
require $A=_{\rho}B$, both paths to be $1$-generic and below $0''$, and the
perfect-family conclusion with the same upper bound $h$.
\end{corollary}
\begin{proof}
If $h$ is bounded by $M$, then $A\triangle B$ has at most $M$ elements.
Finite changes preserve Turing degree, so $A\eqT B$. A noncomputable common
degree contradicts minimal-pairhood.

If $h$ is unbounded, let
$g(n)=\min\{h(n),\lfloor\sqrt n\rfloor\}$.
Both terms are nondecreasing and unbounded, so their minimum is unbounded;
$g$ is computable and sublinear. Apply Theorem~\ref{thm:main} to $g$, and
use $g(n)\le h(n)$.
\end{proof}

The quantifier over $h$ allows bounds much slower than logarithmic growth.
However, it is essential that an unbounded monotone budget eventually grants
one more variation coordinate. A fixed finite number of changes can never
produce these minimal pairs.

\subsection{A perfect family with classical Hausdorff dimension zero}
Equip $\Cantor$ with its usual ultrametric, where a cylinder determined by
$n$ initial bits has diameter at most $2^{-n}$.

\begin{corollary}\label{cor:dimension}
The perfect set $P$ constructed for a sublinear order has classical Hausdorff
dimension zero.
\end{corollary}
\begin{proof}
All its length-$n$ prefixes agree off the fixed set $V\cap[0,n)$.
Thus at most $2^{|V\cap[0,n)|}\le2^{h(n)}$ length-$n$ cylinders meet $P$.
For every real $t>0$, these cylinders provide a cover with total $t$-cost at
most
\[
                         2^{h(n)}(2^{-n})^t=2^{h(n)-tn}\longrightarrow0.
\]
The diameters also tend to zero, so the $t$-dimensional Hausdorff measure is
zero for every $t>0$. The dimension is therefore zero. This is a statement
about the \emph{classical dimension of the family}, not the effective
dimension of each individual path.
\end{proof}

\subsection{No computable mask may reveal infinitely much common data}
The theorem gives a computable \emph{counting bound}, not a computable
superset of the variation set with an infinite computable complement. The
following explains why.

\begin{proposition}[A computable common mask is an obstruction]\label{prop:mask}
If $X$ is $1$-generic, $C\subseteq\N$ is infinite computable, and $Y$ agrees
with $X$ on $C$, then $X,Y$ do not form a Turing minimal pair.
\end{proposition}
\begin{proof}
Let $c_0<c_1<\cdots$ be the computable increasing enumeration of $C$ and put
$Z(n)=X(c_n)=Y(c_n)$. Then $Z\leT X,Y$.
To see that $Z$ is noncomputable, suppose that a computable binary $z$
equals $Z$. The set of strings $\sigma$ for which
$\sigma(c_n)\ne z(n)$ for some $c_n<|\sigma|$ is c.e. and dense: any prefix
can be extended at a later coordinate of $C$ to disagree with $z$.
Genericity forces $X$ to have a prefix in that set, a contradiction.
Thus a noncomputable function is computable from both $X$ and $Y$.
\end{proof}

In particular, the complement of the common variation set in the constructed
family contains no infinite computable subset. Indeed, such a subset would
be a common mask for every pair of distinct members. More generally, no
infinite computable subset can lie in the agreement set of a selected pair.
This does not conflict with that agreement set having density one: high
density does not provide a computable enumeration of usable common positions.

This is the obstruction to the most tempting naive construction: start with
one generic set, alter it on a predetermined computable sparse set, and hope
for a minimal pair. The unchanged computable complement retains noncomputable
information accessible from both oracles. Our construction instead chooses
its sparse positions through the noncomputable stage decisions.

\subsection{What the counterexample does not establish}
The result concerns \emph{least} Turing degrees of representative spectra.
It does not prove that the spectrum has no minimal elements, nor that every
coarse class contains a Turing minimal pair. The theorem is existential in
the coarse class; it does not keep an arbitrary prescribed $1$-generic set
as one of the paths.

Nor does it settle C2 about effective-dense equivalence. An effective-dense
description may omit an answer but may never supply a wrong numerical value.
A density-zero change therefore need not preserve effective-dense descriptions,
and the identity reduction in \eqref{eq:identical-CD} is specific to coarse
descriptions. The report itself records this distinction
\cite[Section~5, C2]{Input}. No effective-dense equivalence is inferred here.

\section{Verification, effectivity, and reproducibility}
\label{sec:verification}

\subsection{The oracle boundary is explicit}
The construction is a $0''$-computable sequence of finite choices. Its
noncomputable decisions have the following precise complexity.
\begin{center}
\begin{tabularx}{\textwidth}{@{}Xll@{}}
\toprule
Decision & Form & Sufficient oracle\\
\midrule
Does a prefix have an extension in a c.e. set $W_e$?
  & $\Sigma^0_1$ & $0'$\\
Does an admissible finite disagreement exist?
  & $\Sigma^0_1$ & $0'$\\
Does a whole partiality cylinder exist?
  & $\Sigma^0_2$ & $0''$\\
Does a proposed cylinder have no finite halting extension?
  & $\Pi^0_1$ & $0'$\\
Is a finite paired suffix admissible?
  & Decidable & None\\
\bottomrule
\end{tabularx}
\end{center}
The common-output program in case C is an ordinary computable search once
the finite prefix and functional index are fixed. There is no claim that
one can uniformly recognize case C without the preceding oracle decisions.

The proof is finite-extension, not a finite-time algorithm for the complete
reals. No halting oracle was accessed during the experiments. In particular,
no finite simulation is offered as evidence that an arbitrary partiality
cylinder really has no future convergence.

\subsection{Executed finite checks}
The accompanying standard-library Python script
\code{checks/verify\_finite\_core.py} was executed successfully. It produces
\code{checks/results.json} and a readable transcript. Its checks are:
\begin{center}
\begin{tabularx}{\textwidth}{@{}Xr@{}}
\toprule
Finite property checked & Number of cases\\
\midrule
Ordered hypercube endpoint pairs, dimensions $0$ through $8$ & 87,381\\
Consecutive cube edges, including common-suffix preservation & 334,962\\
Pairs of total Boolean lookup tables, dimensions $0$ through $3$ & 65,812\\
Pairs of partial lookup tables, dimensions $0$ through $2$ & 6,651\\
Global paired-suffix lifts with every prefix budget checked & 16,380\\
Sparse-coding error-set inclusions and decoding checks & 512\\
Strict-majority block checks & 556\\
\bottomrule
\end{tabularx}
\end{center}
Of the 16,380 suffix lifts, 3,852 were admissible and 12,528 were correctly
rejected. The lookup-table model found only the two constant no-split total
Boolean pairs at each tested dimension. In the partial model, an undefined
table entry \emph{means}, by the definition of that finite model, divergence
on a whole cylinder; it is not a timeout standing in for an unknown answer.

A separate finite structural run carried out six synchronized splitting
levels, ending with $64$ leaves of length $4,098$ and only $12$ global
variation coordinates, while obeying the logarithmic budget at every prefix.
It used fixed seed $20260918$ to choose finite common and one-bit extensions.
It is a test of the budget and lifting mechanics, not a simulation of the
infinite minimal-pair or genericity requirements.

The script also includes a deliberately invalid local-budget shortcut.
On a four-leaf frontier with three global variation positions, one selected
pair differs at only one old position. Charging that pair's count instead of
the global count incorrectly admits a three-bit extension. The regression
check verifies that the correct global test rejects it. This is a concrete
failure mode that a merely pairwise implementation would miss.

\subsection{What remains outside the executable checks}
The infinite theorem rests on the mathematical proof: the density argument
inside Lemma~\ref{lem:trilemma}, the jump-decidable finite-stage choices,
permanent preservation of requirements, and the limit construction.
There is no Lean file claimed to certify these steps, no external referee
report, and no assertion that the Python checks are a substitute for either.
The finite checks make bookkeeping defects easier to detect; they do not
turn the construction into a computable one.

A formalization can isolate four modules with explicit interfaces: finite
strings and finite-use computations; the trilemma; legal-frontier lifting;
and the fair two-path or perfect-family schedule. The nonuniformity of the
case-C program and the distinction between finite divergence and a
partiality cylinder are the two places most deserving a close independent
proof audit.

\section{Conclusion and a precise next refinement}

C1, in the exact formulation printed in the supplied report, has a negative
answer. The known-theorem route establishes this for every $1$-generic set.
The direct route developed here produces two $0''$-computable witnesses and
then a perfect family, with all pairwise Turing information intersections
trivial and all variation confined to an arbitrarily slowly growing
computable counting budget. The main finite mechanism is simple to state:
once one future bit of disagreement is affordable, hypercube connectivity
forces a no-split, no-partiality configuration to have a computable output.

The global budget, not a fixed computable sparse mask, lets this mechanism
coexist with coarse equivalence and individual genericity. The spectrum
calculation further shows that the phenomenon does not depend on choosing
uniform rather than nonuniform coarse classes: those representative spectra
are equal at every fixed base function.

One sharply delimited continuation is to ask whether the \emph{same prescribed
budget theorem} admits two witnesses below $0'$ rather than $0''$.
The present proof does not provide that improvement, because existence of a
partiality cylinder is a $\Sigma^0_2$ decision. This is a refinement of the
construction, not a claim that the question is new or literature-open.
Resolving it would require replacing or approximating that decision while
preserving the global budget and the permanent pair requirements. No
hypothesis about such an improvement is used anywhere in the results above.

\appendix
\section{Oracle-level construction pseudocode}
\label{app:pseudocode}

The following pseudocode is a mathematical specification. The routines whose
names begin with \code{ASK} are oracle queries, not ordinary implemented
subroutines. A witness search is run only after a positive oracle answer.
All frontier extensions preserve its current labels as prefixes.

\begin{lstlisting}
BUFFER(F):
    r := number of global variation coordinates of F
    L := current common length
    find the least Lstar >= L with h(Lstar) >= r + 1
    append 0^(Lstar-L) to every leaf

PAIR_REQUIREMENT(F, i, j, e, d):
    BUFFER(F)
    sigma, tau := current labels of leaves i, j
    r, L := global variation count, common length
    admissible(u,v) := equal lengths and
        r + Hamming(u[:t],v[:t]) <= h(L+t) for every t

    if ASK_0prime(exists admissible finite disagreement):
        find n,u,v and halting traces witnessing disagreement
        append v to leaf j and u to every other leaf
        return

    if ASK_0doubleprime(exists partiality cylinder on either side):
        find a side,n,u using 0prime tests of no halting extension
        append u to every leaf
        return

    # The proved connectivity lemma gives a computable common output.
    # No oracle computation is hidden in that output program.
    return

PERFECT_LEVEL(F_s, s):
    BUFFER(F_s)
    replace each leaf sigma by children sigma0, sigma1
    for each new leaf beta and each e <= s:
        if ASK_0prime(leaf beta has an extension in W_e):
            find such an extension suffix u
            append u to every leaf
        else:
            retain the current prefix as an avoidance witness
    for each ordered pair beta != gamma:
        for each e,d <= s:
            PAIR_REQUIREMENT(F, beta, gamma, e, d)
    freeze the resulting frontier as F_(s+1)
\end{lstlisting}

The iteration order in every finite loop may be fixed lexicographically.
Witnesses may be chosen by the first successful step of a fixed dovetailing.
These conventions make the relative construction completely specified once
standard enumerations of functionals and c.e. sets are fixed. They do not
change any theorem.

\section{Package contents and source provenance}

The package contains the article's \code{.tex} and \code{.pdf}, the
standard-library verification script and its outputs, a build script and
Makefile, a source audit, and an unchanged copy of the supplied
\code{turing\_degrees\_unified.tex}. A SHA-256 manifest records the delivered
files. No third-party font files or copies of the externally cited papers
are bundled.

The source audit distinguishes three roles: the uploaded report specifies
the selected question; Gerdes supplies the exact published-preprint question;
and Hirschfeldt--Jockusch--Kuyper--Schupp supply the older theorem that already
implies its negative answer. The independent construction is not attributed
to any of these sources. The preservation of the input file is for
traceability, not an endorsement of the open-status labels left in that file.

\begin{thebibliography}{9}
\raggedright

\bibitem{Input}
\emph{Turing Degrees: A Unified Research Report}.
User-supplied LaTeX source \code{turing\_degrees\_unified.tex}, unified edition
18 September 2026. Section~5: ``Coarse and effective-dense representative
spectra,'' especially C1 and the sparse-coding and minimal-pair diagnostics.
An unchanged copy is included in \code{input/}.

\bibitem{Gerdes}
Peter M. Gerdes.
\emph{Comparing Notions of Dense Computability on $\omega^\omega$ and $2^\omega$}.
arXiv:2508.06925v1, 9 August 2025.
Question~7, printed p.~64.
\url{https://arxiv.org/abs/2508.06925}.
Accessed 18 September 2026.

\bibitem{HJKS}
Denis R. Hirschfeldt, Carl G. Jockusch, Jr., Rutger Kuyper, and Paul E. Schupp.
\emph{Coarse reducibility and algorithmic randomness}.
Journal of Symbolic Logic \textbf{81} (2016), 1028--1046.
DOI: \href{https://doi.org/10.1017/jsl.2015.70}{10.1017/jsl.2015.70}.
Preprint arXiv:1505.01707 (2015), especially Definition~3.1 and Theorem~4.2
(printed p.~16 of the inspected preprint).
\url{https://arxiv.org/abs/1505.01707}.
Publication metadata checked against the author's publication page:
\url{https://www.math.uchicago.edu/~drh/Papers/coarsereducibility.html}.
Accessed 18 September 2026.

\end{thebibliography}
\end{document}
